\section{Results: Adaptation Budget and Drift}
\label{sec:res_budget}

\begin{figure}[t]
  \centering
  \includegraphics[width=0.94\linewidth]{figures/budget_curves.png}
  \caption{Adapters trained for a fixed 6000 steps on nested subsets (log axis), zero-shot dashed. Left: in-domain Dice. Centre: far-OOD Dice, mean of BUSI and CBIS-DDSM. Right: far-OOD representation drift ($1 - \mathrm{CKA}$ to the base model) of the decoder's upscaled embedding (circles) and of the encoder neck output (squares, dashed); decoder-only fine-tuning leaves the encoder untouched. Dice: mean and one std over three seeds at 50, 250 and 1000 images; the 2595-image point and the drift panel are seed 0.}
  \label{fig:budget}
\end{figure}

Figure~\ref{fig:budget} shows how the collapse depends on the budget and what moves inside the model when it happens. Decoder-only fine-tuning and encoder-only LoRA improve on both axes with every increase in budget and stay above the zero-shot far-OOD reference of 0.758 in every seed (0.840 to 0.861 and 0.811 to 0.864). LoRA with a trainable decoder is the only method that crosses it: already 0.02 to 0.05 below the others at 50 images ($0.790 \pm 0.027$, above zero-shot in each seed), it falls to $0.670 \pm 0.031$ at 250, below zero-shot in all three seeds, recovers to $0.774 \pm 0.019$ at 1000, above it in all three, and ends at 0.638 with all 2595 images, while its in-domain Dice rises monotonically. The drift panel says why. LoRA's far-OOD decoder drift triples between 50 and 250 images, from 0.12 to 0.36, exactly where its Dice collapses, and stays there. Decoder-only fine-tuning retrains every decoder weight yet keeps far-OOD decoder drift below 0.06 at all budgets; encoder-only LoRA moves the encoder twice as much as LoRA (encoder drift 0.25 against 0.12) and keeps far-OOD decoder drift below 0.09; in-domain decoder drift is similar for all three (0.07 to 0.15). The collapse is thus not a symptom of too little data (every budget receives 6000 steps), of decoder training as such, or of encoder movement as such: it appears when encoder and decoder adapt jointly without an anchor, as far-OOD decoder drift that more data does not undo. The dip at 250 images and the partial recovery at 1000 are reproduced in every seed; the recovery occurs at unchanged far-OOD decoder drift (0.35), so drift marks the onset of the collapse rather than tracking its depth.
